
        You are a research assistant AI tasked with generating a scientific paper based on provided literature. Follow these steps:
    1. Analyze the given References.
    2. Identify gaps in existing research to establish the motivation for a new study.
    3. Propose a main idea for a new research work.
    4. Write the paper's main content in LaTeX format, including all the following parts, please must have tables:
    - Title
    - Abstract
    - Introduction
    - Related Work
    - Methods/Experiment
    - Results with tables
    - Discussion
    - Conclusion
    - Contributions
    Ensure all content is original, academically rigorous, and follows standard scientific writing conventions.

        Here are the references to be used for generating the paper.
        You MUST base the paper on these references and integrate them naturally.

        References:
        --------------------
        @misc{narasimhappa2026hybridlstmukfframeworkankle,
      title={Hybrid LSTM-UKF Framework: Ankle Angle and Ground Reaction Force Estimation}, 
      author={Mundla Narasimhappa and Praveen Kumar},
      year={2026},
      eprint={2601.06473},
      archivePrefix={arXiv},
      primaryClass={eess.SY},
      url={https://arxiv.org/abs/2601.06473},
      abstract = {Accurate prediction of joint kinematics and kinetics is essential for advancing gait analysis and developing intelligent assistive systems such as prosthetics and exoskeletons. This study presents a hybrid LSTM-UKF framework for estimating ankle angle and ground reaction force (GRF) across varying walking speeds. A multimodal sensor fusion strategy integrates force plate data, knee angle, and GRF signals to enrich biomechanical context. Model performance was evaluated using RMSE and $R^2$ under subject-specific validation. The LSTM-UKF consistently outperformed standalone LSTM and UKF models, achieving up to 18.6\\\% lower RMSE for GRF prediction at 3 km/h. Additionally, UKF integration improved robustness, reducing ankle angle RMSE by up to 22.4\\\% compared to UKF alone at 1 km/h. These results underscore the effectiveness of hybrid architectures for reliable gait prediction across subjects and walking conditions.}
}@misc{ma2026challengesresearchdirectionslarge,
      title={Challenges and Research Directions for Large Language Model Inference Hardware}, 
      author={Xiaoyu Ma and David Patterson},
      year={2026},
      eprint={2601.05047},
      archivePrefix={arXiv},
      primaryClass={cs.AR},
      doi={https://doi.org/10.1109/MC.2026.3652916},
      url={https://arxiv.org/abs/2601.05047},
      abstract = {Large Language Model (LLM) inference is hard. The autoregressive Decode phase of the underlying Transformer model makes LLM inference fundamentally different from training. Exacerbated by recent AI trends, the primary challenges are memory and interconnect rather than compute. To address these challenges, we highlight four architecture research opportunities: High Bandwidth Flash for 10X memory capacity with HBM-like bandwidth; Processing-Near-Memory and 3D memory-logic stacking for high memory bandwidth; and low-latency interconnect to speedup communication. While our focus is datacenter AI, we also review their applicability for mobile devices.}
}@misc{spaan2026reducingcomputewastellms,
      title={Reducing Compute Waste in LLMs through Kernel-Level DVFS}, 
      author={Jeffrey Spaan and Kuan-Hsun Chen and Ana-Lucia Varbanescu},
      year={2026},
      eprint={2601.08539},
      archivePrefix={arXiv},
      primaryClass={cs.PF},
      url={https://arxiv.org/abs/2601.08539},
      abstract = {The rapid growth of AI has fueled the expansion of accelerator- or GPU-based data centers. However, the rising operational energy consumption has emerged as a critical bottleneck and a major sustainability concern. Dynamic Voltage and Frequency Scaling (DVFS) is a well-known technique used to reduce energy consumption, and thus improve energy-efficiency, since it requires little effort and works with existing hardware. Reducing the energy consumption of training and inference of Large Language Models (LLMs) through DVFS or power capping is feasible: related work has shown energy savings can be significant, but at the cost of significant slowdowns. In this work, we focus on reducing waste in LLM operations: i.e., reducing energy consumption without losing performance. We propose a fine-grained, kernel-level, DVFS approach that explores new frequency configurations, and prove these save more energy than previous, pass- or iteration-level solutions. For example, for a GPT-3 training run, a pass-level approach could reduce energy consumption by 2\% (without losing performance), while our kernel-level approach saves as much as 14.6\% (with a 0.6\% slowdown). We further investigate the effect of data and tensor parallelism, and show our discovered clock frequencies translate well for both. We conclude that kernel-level DVFS is a suitable technique to reduce waste in LLM operations, providing significant energy savings with negligible slow-down.}
}@misc{bakhmach2026roofsrobustbiomarkerfeature,
      title={ROOFS: RObust biOmarker Feature Selection}, 
      author={Anastasiia Bakhmach and Paul Dufossé and Andrea Vaglio and Florence Monville and Laurent Greillier and Fabrice Barlési and Sébastien Benzekry},
      year={2026},
      eprint={2601.05151},
      archivePrefix={arXiv},
      primaryClass={stat.ML},
      url={https://arxiv.org/abs/2601.05151},
      abstract = {Feature selection (FS) is essential for biomarker discovery and in the analysis of biomedical datasets. However, challenges such as high-dimensional feature space, low sample size, multicollinearity, and missing values make FS non-trivial. Moreover, FS performances vary across datasets and predictive tasks. We propose roofs, a Python package available at https://gitlab.inria.fr/compo/roofs, designed to help researchers in the choice of FS method adapted to their problem. Roofs benchmarks multiple FS methods on the user's data and generates reports that summarize a comprehensive set of evaluation metrics, including downstream predictive performance estimated using optimism correction, stability, reliability of individual features, and true positive and false positive rates assessed on semi-synthetic data with a simulated outcome. We demonstrate the utility of roofs on data from the PIONeeR clinical trial, aimed at identifying predictors of resistance to anti-PD-(L)1 immunotherapy in lung cancer. The PIONeeR dataset contained 374 multi-source blood and tumor biomarkers from 435 patients. A reduced subset of 214 features was obtained through iterative variance inflation factor pre-filtering. Of the 34 FS methods gathered in roofs, we evaluated 23 in combination with 11 classifiers (253 models in total) and identified a filter based on the union of Benjamini-Hochberg false discovery rate-adjusted p-values from t-test and logistic regression as the optimal approach, outperforming other methods including the widely used LASSO. We conclude that comprehensive benchmarking with roofs has the potential to improve the robustness and reproducibility of FS discoveries and increase the translational value of clinical models.}
}@misc{ozgur2026vitntfiqatrainingfreefaceimage,
      title={ViTNT-FIQA: Training-Free Face Image Quality Assessment with Vision Transformers}, 
      author={Guray Ozgur and Eduarda Caldeira and Tahar Chettaoui and Jan Niklas Kolf and Marco Huber and Naser Damer and Fadi Boutros},
      year={2026},
      eprint={2601.05741},
      archivePrefix={arXiv},
      primaryClass={cs.CV},
      url={https://arxiv.org/abs/2601.05741},
      abstract = {Face Image Quality Assessment (FIQA) is essential for reliable face recognition systems. Current approaches primarily exploit only final-layer representations, while training-free methods require multiple forward passes or backpropagation. We propose ViTNT-FIQA, a training-free approach that measures the stability of patch embedding evolution across intermediate Vision Transformer (ViT) blocks. We demonstrate that high-quality face images exhibit stable feature refinement trajectories across blocks, while degraded images show erratic transformations. Our method computes Euclidean distances between L2-normalized patch embeddings from consecutive transformer blocks and aggregates them into image-level quality scores. We empirically validate this correlation on a quality-labeled synthetic dataset with controlled degradation levels. Unlike existing training-free approaches, ViTNT-FIQA requires only a single forward pass without backpropagation or architectural modifications. Through extensive evaluation on eight benchmarks (LFW, AgeDB-30, CFP-FP, CALFW, Adience, CPLFW, XQLFW, IJB-C), we show that ViTNT-FIQA achieves competitive performance with state-of-the-art methods while maintaining computational efficiency and immediate applicability to any pre-trained ViT-based face recognition model.}
}@misc{zhang2026weightscodeextractinginterpretable,
      title={Weights to Code: Extracting Interpretable Algorithms from the Discrete Transformer}, 
      author={Yifan Zhang and Wei Bi and Kechi Zhang and Dongming Jin and Jie Fu and Zhi Jin},
      year={2026},
      eprint={2601.05770},
      archivePrefix={arXiv},
      primaryClass={cs.LG},
      url={https://arxiv.org/abs/2601.05770},
      abstract = {Algorithm extraction aims to synthesize executable programs directly from models trained on specific algorithmic tasks, enabling de novo algorithm discovery without relying on human-written code. However, extending this paradigm to Transformer is hindered by superposition, where entangled features encoded in overlapping directions obstruct the extraction of symbolic expressions. In this work, we propose the Discrete Transformer, an architecture explicitly engineered to bridge the gap between continuous representations and discrete symbolic logic. By enforcing a strict functional disentanglement, which constrains Numerical Attention to information routing and Numerical MLP to element-wise arithmetic, and employing temperature-annealed sampling, our method effectively facilitates the extraction of human-readable programs. Empirically, the Discrete Transformer not only achieves performance comparable to RNN-based baselines but crucially extends interpretability to continuous variable domains. Moreover, our analysis of the annealing process shows that the efficient discrete search undergoes a clear phase transition from exploration to exploitation. We further demonstrate that our method enables fine-grained control over synthesized programs by imposing inductive biases. Collectively, these findings establish the Discrete Transformer as a robust framework for demonstration-free algorithm discovery, offering a rigorous pathway toward Transformer interpretability.}
}@misc{fontanari2026hierarchicalpoolingexplainabilitygraph,
      title={Hierarchical Pooling and Explainability in Graph Neural Networks for Tumor and Tissue-of-Origin Classification Using RNA-seq Data}, 
      author={Thomas Vaitses Fontanari and Mariana Recamonde-Mendoza},
      year={2026},
      eprint={2601.06381},
      archivePrefix={arXiv},
      primaryClass={cs.LG},
      url={https://arxiv.org/abs/2601.06381},
      abstract = {This study explores the use of graph neural networks (GNNs) with hierarchical pooling and multiple convolution layers for cancer classification based on RNA-seq data. We combine gene expression data from The Cancer Genome Atlas (TCGA) with a precomputed STRING protein-protein interaction network to classify tissue origin and distinguish between normal and tumor samples. The model employs Chebyshev graph convolutions (K=2) and weighted pooling layers, aggregating gene clusters into'supernodes' across multiple coarsening levels. This approach enables dimensionality reduction while preserving meaningful interactions. Saliency methods were applied to interpret the model by identifying key genes and biological processes relevant to cancer. Our findings reveal that increasing the number of convolution and pooling layers did not enhance classification performance. The highest F1-macro score (0.978) was achieved with a single pooling layer. However, adding more layers resulted in over-smoothing and performance degradation. However, the model proved highly interpretable through gradient methods, identifying known cancer-related genes and highlighting enriched biological processes, and its hierarchical structure can be used to develop new explainable architectures. Overall, while deeper GNN architectures did not improve performance, the hierarchical pooling structure provided valuable insights into tumor biology, making GNNs a promising tool for cancer biomarker discovery and interpretation}
}@misc{yasin2026showdonttell,
      title={Show, don't tell -- Providing Visual Error Feedback for Handwritten Documents}, 
      author={Said Yasin and Torsten Zesch},
      year={2026},
      eprint={2601.09586},
      archivePrefix={arXiv},
      primaryClass={cs.CV},
      url={https://arxiv.org/abs/2601.09586},
      abstract = {Handwriting remains an essential skill, particularly in education. Therefore, providing visual feedback on handwritten documents is an important but understudied area. We outline the many challenges when going from an image of handwritten input to correctly placed informative error feedback. We empirically compare modular and end-to-end systems and find that both approaches currently do not achieve acceptable overall quality. We identify the major challenges and outline an agenda for future research.}
}@misc{chen2026xlinearlightweightaccuratemlpbased,
      title={XLinear: A Lightweight and Accurate MLP-Based Model for Long-Term Time Series Forecasting with Exogenous Inputs}, 
      author={Xinyang Chen and Huidong Jin and Yu Huang and Zaiwen Feng},
      year={2026},
      eprint={2601.09237},
      archivePrefix={arXiv},
      primaryClass={cs.LG},
      url={https://arxiv.org/abs/2601.09237},
      abstract = {Despite the prevalent assumption of uniform variable importance in long-term time series forecasting models, real world applications often exhibit asymmetric causal relationships and varying data acquisition costs. Specifically, cost-effective exogenous data (e.g., local weather) can unilaterally influence dynamics of endogenous variables, such as lake surface temperature. Exploiting these links enables more effective forecasts when exogenous inputs are readily available. Transformer-based models capture long-range dependencies but incur high computation and suffer from permutation invariance. Patch-based variants improve efficiency yet can miss local temporal patterns. To efficiently exploit informative signals across both the temporal dimension and relevant exogenous variables, this study proposes XLinear, a lightweight time series forecasting model built upon MultiLayer Perceptrons (MLPs). XLinear uses a global token derived from an endogenous variable as a pivotal hub for interacting with exogenous variables, and employs MLPs with sigmoid activation to extract both temporal patterns and variate-wise dependencies. Its prediction head then integrates these signals to forecast the endogenous series. We evaluate XLinear on seven standard benchmarks and five real-world datasets with exogenous inputs. Compared with state-of-the-art models, XLinear delivers superior accuracy and efficiency for both multivariate forecasts and univariate forecasts influenced by exogenous inputs.}
}@misc{fang2026integratedcrossurbanaccidentprevention,
      title={Toward an Integrated Cross-Urban Accident Prevention System: A Multi-Task Spatial-Temporal Learning Framework for Urban Safety Management}, 
      author={Jiayu Fang and Zhiqi Shao and Haoning Xi and Boris Choy and Junbin Gao},
      year={2026},
      eprint={2601.05521},
      archivePrefix={arXiv},
      primaryClass={cs.LG},
      url={https://arxiv.org/abs/2601.05521},
      abstract = {The development of a cross-city accident prevention system is particularly challenging due to the heterogeneity, inconsistent reporting, and inherently clustered, sparse, cyclical, and noisy nature of urban accident data. These intrinsic data properties, combined with fragmented governance and incompatible reporting standards, have long hindered the creation of an integrated, cross-city accident prevention framework. To address this gap, we propose the Mamba Local-ttention Spatial-Temporal Network MLA-STNet, a unified system that formulates accident risk prediction as a multi-task learning problem across multiple cities. MLA-STNet integrates two complementary modules: (i)the Spatio-Temporal Geographical Mamba-Attention (STG-MA), which suppresses unstable spatio-temporal fluctuations and strengthens long-range temporal dependencies; and (ii) the Spatio-Temporal Semantic Mamba-Attention (STS-MA), which mitigates cross-city heterogeneity through a shared-parameter design that jointly trains all cities while preserving individual semantic representation spaces. We validate the proposed framework through 75 experiments under two forecasting scenarios, full-day and high-frequency accident periods, using real-world datasets from New York City and Chicago. Compared with the state-of-the-art baselines, MLA-STNet achieves up to 6\% lower RMSE, 8\% higher Recall, and 5\% higher MAP, while maintaining less than 1\% performance variation under 50\% input noise. These results demonstrate that MLA-STNet effectively unifies heterogeneous urban datasets within a scalable, robust, and interpretable Cross-City Accident Prevention System, paving the way for coordinated and data-driven urban safety management.}
}@misc{tan2026kernelalignmentbasedmultiviewunsupervised,
      title={Kernel Alignment-based Multi-view Unsupervised Feature Selection with Sample-level Adaptive Graph Learning}, 
      author={Yalan Tan and Yanyong Huang and Zongxin Shen and Dongjie Wang and Fengmao Lv and Tianrui Li},
      year={2026},
      eprint={2601.07288},
      archivePrefix={arXiv},
      primaryClass={cs.LG},
      url={https://arxiv.org/abs/2601.07288},
      abstract = {Although multi-view unsupervised feature selection (MUFS) has demonstrated success in dimensionality reduction for unlabeled multi-view data, most existing methods reduce feature redundancy by focusing on linear correlations among features but often overlook complex nonlinear dependencies. This limits the effectiveness of feature selection. In addition, existing methods fuse similarity graphs from multiple views by employing sample-invariant weights to preserve local structure. However, this process fails to account for differences in local neighborhood clarity among samples within each view, thereby hindering accurate characterization of the intrinsic local structure of the data. In this paper, we propose a Kernel Alignment-based multi-view unsupervised FeatUre selection with Sample-level adaptive graph lEarning method (KAFUSE) to address these issues. Specifically, we first employ kernel alignment with an orthogonal constraint to reduce feature redundancy in both linear and nonlinear relationships. Then, a cross-view consistent similarity graph is learned by applying sample-level fusion to each slice of a tensor formed by stacking similarity graphs from different views, which automatically adjusts the view weights for each sample during fusion. These two steps are integrated into a unified model for feature selection, enabling mutual enhancement between them. Extensive experiments on real multi-view datasets demonstrate the superiority of KAFUSE over state-of-the-art methods.}
}@misc{zhao2026fastefficienteffectivelonghorizon,
      title={FaST: Efficient and Effective Long-Horizon Forecasting for Large-Scale Spatial-Temporal Graphs via Mixture-of-Experts}, 
      author={Yiji Zhao and Zihao Zhong and Ao Wang and Haomin Wen and Ming Jin and Yuxuan Liang and Huaiyu Wan and Hao Wu},
      year={2026},
      eprint={2601.05174},
      archivePrefix={arXiv},
      primaryClass={cs.LG},
      doi={https://doi.org/10.1145/3770854.3780165},
      url={https://arxiv.org/abs/2601.05174},
      abstract = {Spatial-Temporal Graph (STG) forecasting on large-scale networks has garnered significant attention. However, existing models predominantly focus on short-horizon predictions and suffer from notorious computational costs and memory consumption when scaling to long-horizon predictions and large graphs. Targeting the above challenges, we present FaST, an effective and efficient framework based on heterogeneity-aware Mixture-of-Experts (MoEs) for long-horizon and large-scale STG forecasting, which unlocks one-week-ahead (672 steps at a 15-minute granularity) prediction with thousands of nodes. FaST is underpinned by two key innovations. First, an adaptive graph agent attention mechanism is proposed to alleviate the computational burden inherent in conventional graph convolution and self-attention modules when applied to large-scale graphs. Second, we propose a new parallel MoE module that replaces traditional feed-forward networks with Gated Linear Units (GLUs), enabling an efficient and scalable parallel structure. Extensive experiments on real-world datasets demonstrate that FaST not only delivers superior long-horizon predictive accuracy but also achieves remarkable computational efficiency compared to state-of-the-art baselines. Our source code is available at: https://github.com/yijizhao/FaST.}
}@misc{yin2026latencyprismonlinenonintrusivelatency,
      title={LatencyPrism: Online Non-intrusive Latency Sculpting for SLO-Guaranteed LLM Inference}, 
      author={Du Yin and Jiayi Ren and Xiayu Sun and Tianyao Zhou and Haizhu Zhou and Ruiyan Ma and Danyang Zhang},
      year={2026},
      eprint={2601.09258},
      archivePrefix={arXiv},
      primaryClass={cs.DC},
      url={https://arxiv.org/abs/2601.09258},
      abstract = {LLM inference latency critically determines user experience and operational costs, directly impacting throughput under SLO constraints. Even brief latency spikes degrade service quality despite acceptable average performance. However, distributed inference environments featuring diverse software frameworks and XPU architectures combined with dynamic workloads make latency analysis challenging. Constrained by intrusive designs that necessitate service restarts or even suspension, and by hardware-bound implementations that fail to adapt to heterogeneous inference environments, existing AI profiling methods are often inadequate for real-time production analysis.   We present LatencyPrism, the first zero-intrusion multi-platform latency sculpting system. It aims to break down the inference latency across pipeline, proactively alert on inference latency anomalies, and guarantee adherence to SLOs, all without requiring code modifications or service restarts. LatencyPrism has been deployed across thousands of XPUs for over six months. It enables low-overhead real-time monitoring at batch level with alerts triggered in milliseconds. This approach distinguishes between workload-driven latency variations and anomalies indicating underlying issues with an F1-score of 0.98. We also conduct extensive experiments and investigations into root cause analysis to demonstrate LatencyPrism's capability.}
}@misc{shin2026agdcautoregressivegenerationvariablelength,
      title={AGDC: Autoregressive Generation of Variable-Length Sequences with Joint Discrete and Continuous Spaces}, 
      author={Yeonsang Shin and Insoo Kim and Bongkeun Kim and Keonwoo Bae and Bohyung Han},
      year={2026},
      eprint={2601.05680},
      archivePrefix={arXiv},
      primaryClass={cs.LG},
      url={https://arxiv.org/abs/2601.05680},
      abstract = {Transformer-based autoregressive models excel in data generation but are inherently constrained by their reliance on discretized tokens, which limits their ability to represent continuous values with high precision. We analyze the scalability limitations of existing discretization-based approaches for generating hybrid discrete-continuous sequences, particularly in high-precision domains such as semiconductor circuit designs, where precision loss can lead to functional failure. To address the challenge, we propose AGDC, a novel unified framework that jointly models discrete and continuous values for variable-length sequences. AGDC employs a hybrid approach that combines categorical prediction for discrete values with diffusion-based modeling for continuous values, incorporating two key technical components: an end-of-sequence (EOS) logit adjustment mechanism that uses an MLP to dynamically adjust EOS token logits based on sequence context, and a length regularization term integrated into the loss function. Additionally, we present ContLayNet, a large-scale benchmark comprising 334K high-precision semiconductor layout samples with specialized evaluation metrics that capture functional correctness where precision errors significantly impact performance. Experiments on semiconductor layouts (ContLayNet), graphic layouts, and SVGs demonstrate AGDC's superior performance in generating high-fidelity hybrid vector representations compared to discretization-based and fixed-schema baselines, achieving scalable high-precision generation across diverse domains.}
}@misc{hummel2026decodablestructuredlinearprobing,
      title={Decodable but not structured: linear probing enables Underwater Acoustic Target Recognition with pretrained audio embeddings}, 
      author={Hilde I. Hummel and Sandjai Bhulai and Rob D. van der Mei and Burooj Ghani},
      year={2026},
      eprint={2601.08358},
      archivePrefix={arXiv},
      primaryClass={cs.LG},
      url={https://arxiv.org/abs/2601.08358},
      abstract = {Increasing levels of anthropogenic noise from ships contribute significantly to underwater sound pollution, posing risks to marine ecosystems. This makes monitoring crucial to understand and quantify the impact of the ship radiated noise. Passive Acoustic Monitoring (PAM) systems are widely deployed for this purpose, generating years of underwater recordings across diverse soundscapes. Manual analysis of such large-scale data is impractical, motivating the need for automated approaches based on machine learning. Recent advances in automatic Underwater Acoustic Target Recognition (UATR) have largely relied on supervised learning, which is constrained by the scarcity of labeled data. Transfer Learning (TL) offers a promising alternative to mitigate this limitation. In this work, we conduct the first empirical comparative study of transfer learning for UATR, evaluating multiple pretrained audio models originating from diverse audio domains. The pretrained model weights are frozen, and the resulting embeddings are analyzed through classification, clustering, and similarity-based evaluations. The analysis shows that the geometrical structure of the embedding space is largely dominated by recording-specific characteristics. However, a simple linear probe can effectively suppress this recording-specific information and isolate ship-type features from these embeddings. As a result, linear probing enables effective automatic UATR using pretrained audio models at low computational cost, significantly reducing the need for a large amounts of high-quality labeled ship recordings.}
}@misc{taneja2026gpuacceleratedint8quantizationkv,
      title={GPU-Accelerated INT8 Quantization for KV Cache Compression in Large Language Models}, 
      author={Maanas Taneja and Purab Shingvi},
      year={2026},
      eprint={2601.04719},
      archivePrefix={arXiv},
      primaryClass={cs.LG},
      url={https://arxiv.org/abs/2601.04719},
      abstract = {The key-value (KV) cache in large language models presents a significant memory bottleneck during inference, growing linearly with sequence length and often exceeding the memory footprint of model weights themselves. We implement and evaluate GPU-accelerated INT8 quantization for KV cache compression, achieving 4$\\times$ memory reduction with minimal accuracy degradation. We develop four CUDA kernel variants -- naive, tiled, coarsened, and vectorized -- and benchmark them across realistic workload sizes up to 1 billion elements. Our vectorized kernel achieves up to 1,694$\\times$ speedup over CPU baselines while maintaining reconstruction error below 0.004 and attention score error below 0.1 even for 8K-dimensional heads. These results demonstrate that INT8 quantization provides a practical approach for reducing memory pressure in LLM inference with negligible computational overhead (6--58ms) and minimal impact on downstream model behavior}
}@misc{kleiner2026illuminationangularspectrumencoding,
      title={Illumination Angular Spectrum Encoding for Controlling the Functionality of Diffractive Networks}, 
      author={Matan Kleiner and Lior Michaeli and Tomer Michaeli},
      year={2026},
      eprint={2601.04825},
      archivePrefix={arXiv},
      primaryClass={physics.optics},
      url={https://arxiv.org/abs/2601.04825},
      abstract = {Diffractive neural networks have recently emerged as a promising framework for all-optical computing. However, these networks are typically trained for a single task, limiting their potential adoption in systems requiring multiple functionalities. Existing approaches to achieving multi-task functionality either modify the mechanical configuration of the network per task or use a different illumination wavelength or polarization state for each task. In this work, we propose a new control mechanism, which is based on the illumination's angular spectrum. Specifically, we shape the illumination using an amplitude mask that selectively controls its angular spectrum. We employ different illumination masks for achieving different network functionalities, so that the mask serves as a unique task encoder. Interestingly, we show that effective control can be achieved over a very narrow angular range, within the paraxial regime. We numerically illustrate the proposed approach by training a single diffractive network to perform multiple image-to-image translation tasks. In particular, we demonstrate translating handwritten digits into typeset digits of different values, and translating handwritten English letters into typeset numbers and typeset Greek letters, where the type of the output is determined by the illumination's angular components. As we show, the proposed framework can work under different coherence conditions, and can be combined with existing control strategies, such as different wavelengths. Our results establish the illumination angular spectrum as a powerful degree of freedom for controlling diffractive networks, enabling a scalable and versatile framework for multi-task all-optical computing.}
}@misc{chen2026interpretabilityindividualitykneemri,
      title={Interpretability and Individuality in Knee MRI: Patient-Specific Radiomic Fingerprint with Reconstructed Healthy Personas}, 
      author={Yaxi Chen and Simin Ni and Shuai Li and Shaheer U. Saeed and Aleksandra Ivanova and Rikin Hargunani and Jie Huang and Chaozong Liu and Yipeng Hu},
      year={2026},
      eprint={2601.08604},
      archivePrefix={arXiv},
      primaryClass={cs.CV},
      url={https://arxiv.org/abs/2601.08604},
      abstract = {For automated assessment of knee MRI scans, both accuracy and interpretability are essential for clinical use and adoption. Traditional radiomics rely on predefined features chosen at the population level; while more interpretable, they are often too restrictive to capture patient-specific variability and can underperform end-to-end deep learning (DL). To address this, we propose two complementary strategies that bring individuality and interpretability: radiomic fingerprints and healthy personas. First, a radiomic fingerprint is a dynamically constructed, patient-specific feature set derived from MRI. Instead of applying a uniform population-level signature, our model predicts feature relevance from a pool of candidate features and selects only those most predictive for each patient, while maintaining feature-level interpretability. This fingerprint can be viewed as a latent-variable model of feature usage, where an image-conditioned predictor estimates usage probabilities and a transparent logistic regression with global coefficients performs classification. Second, a healthy persona synthesises a pathology-free baseline for each patient using a diffusion model trained to reconstruct healthy knee MRIs. Comparing features extracted from pathological images against their personas highlights deviations from normal anatomy, enabling intuitive, case-specific explanations of disease manifestations. We systematically compare fingerprints, personas, and their combination across three clinical tasks. Experimental results show that both approaches yield performance comparable to or surpassing state-of-the-art DL models, while supporting interpretability at multiple levels. Case studies further illustrate how these perspectives facilitate human-explainable biomarker discovery and pathology localisation.}
}@misc{knoop2026privatellminferenceconsumer,
      title={Private LLM Inference on Consumer Blackwell GPUs: A Practical Guide for Cost-Effective Local Deployment in SMEs}, 
      author={Jonathan Knoop and Hendrik Holtmann},
      year={2026},
      eprint={2601.09527},
      archivePrefix={arXiv},
      primaryClass={cs.LG},
      url={https://arxiv.org/abs/2601.09527},
      abstract = {SMEs increasingly seek alternatives to cloud LLM APIs, which raise data privacy concerns. Dedicated cloud GPU instances offer improved privacy but with limited guarantees and ongoing costs, while professional on-premise hardware (A100, H100) remains prohibitively expensive. We present a systematic evaluation of NVIDIA's Blackwell consumer GPUs (RTX 5060 Ti, 5070 Ti, 5090) for production LLM inference, benchmarking four open-weight models (Qwen3-8B, Gemma3-12B, Gemma3-27B, GPT-OSS-20B) across 79 configurations spanning quantization formats (BF16, W4A16, NVFP4, MXFP4), context lengths (8k-64k), and three workloads: RAG, multi-LoRA agentic serving, and high-concurrency APIs. The RTX 5090 delivers 3.5-4.6x higher throughput than the 5060 Ti with 21x lower latency for RAG, but budget GPUs achieve the highest throughput-per-dollar for API workloads with sub-second latency. NVFP4 quantization provides 1.6x throughput over BF16 with 41\% energy reduction and only 2-4\% quality loss. Self-hosted inference costs $0.001-0.04 per million tokens (electricity only), which is 40-200x cheaper than budget-tier cloud APIs, with hardware breaking even in under four months at moderate volume (30M tokens/day). Our results show that consumer GPUs can reliably replace cloud inference for most SME workloads, except latency-critical long-context RAG, where high-end GPUs remain essential. We provide deployment guidance and release all benchmark data for reproducible SME-scale deployments.}
}@misc{lam2026energyentropyregularizationtruepower,
      title={Energy-Entropy Regularization: The True Power of Minimal Looped Transformers}, 
      author={Wai-Lun Lam},
      year={2026},
      eprint={2601.09588},
      archivePrefix={arXiv},
      primaryClass={cs.LG},
      url={https://arxiv.org/abs/2601.09588},
      abstract = {Recent research suggests that looped Transformers have superior reasoning capabilities compared to standard deep architectures. Current approaches to training single-head looped architectures on benchmark tasks frequently fail or yield suboptimal performance due to a highly non-convex and irregular loss landscape. In these settings, optimization often stagnates in poor local minima and saddle points of the loss landscape, preventing the model from discovering the global minimum point. The internal mechanisms of these single-head looped transformer models remain poorly understood, and training them from scratch remains a significant challenge. In this paper, we propose a novel training framework that leverages Tsallis entropy and Hamiltonian dynamics to transform the geometry of the loss landscape. By treating the parameter updates as a physical flow, we successfully trained a single-head looped Transformer with model dimension $d = 8$ to solve induction head task with input sequence length of 1000 tokens. This success reveals the internal mechanism behind the superior reasoning capability.}
}@misc{jamil2026explainabledeepradiogenomicmolecular,
      title={Explainable Deep Radiogenomic Molecular Imaging for MGMT Methylation Prediction in Glioblastoma}, 
      author={Hasan M Jamil},
      year={2026},
      eprint={2601.07035},
      archivePrefix={arXiv},
      primaryClass={cs.LG},
      url={https://arxiv.org/abs/2601.07035},
      abstract = {Glioblastoma (GBM) is a highly aggressive primary brain tumor with limited therapeutic options and poor prognosis. The methylation status of the O6-methylguanine-DNA methyltransferase (MGMT) gene promoter is a critical molecular biomarker that influences patient response to temozolomide chemotherapy. Traditional methods for determining MGMT status rely on invasive biopsies and are limited by intratumoral heterogeneity and procedural risks. This study presents a radiogenomic molecular imaging analysis framework for the non-invasive prediction of MGMT promoter methylation using multi-parametric magnetic resonance imaging (mpMRI).   Our approach integrates radiomics, deep learning, and explainable artificial intelligence (XAI) to analyze MRI-derived imaging phenotypes and correlate them with molecular labels. Radiomic features are extracted from FLAIR, T1-weighted, T1-contrast-enhanced, and T2-weighted MRI sequences, while a 3D convolutional neural network learns deep representations from the same modalities. These complementary features are fused using both early fusion and attention-based strategies and classified to predict MGMT methylation status.   To enhance clinical interpretability, we apply XAI methods such as Grad-CAM and SHAP to visualize and explain model decisions. The proposed framework is trained on the RSNA-MICCAI Radiogenomic Classification dataset and externally validated on the BraTS 2021 dataset. This work advances the field of molecular imaging by demonstrating the potential of AI-driven radiogenomics for precision oncology, supporting non-invasive, accurate, and interpretable prediction of clinically actionable molecular biomarkers in GBM.}
}@misc{gupta2026performancedeeplearningbasedsegmentation,
      title={Performance of a Deep Learning-Based Segmentation Model for Pancreatic Tumors on Public Endoscopic Ultrasound Datasets}, 
      author={Pankaj Gupta and Priya Mudgil and Niharika Dutta and Kartik Bose and Nitish Kumar and Anupam Kumar and Jimil Shah and Vaneet Jearth and Jayanta Samanta and Vishal Sharma and Harshal Mandavdhare and Surinder Rana and Saroj K Sinha and Usha Dutta},
      year={2026},
      eprint={2601.05937},
      archivePrefix={arXiv},
      primaryClass={cs.CV},
      url={https://arxiv.org/abs/2601.05937},
      abstract = {Background: Pancreatic cancer is one of the most aggressive cancers, with poor survival rates. Endoscopic ultrasound (EUS) is a key diagnostic modality, but its effectiveness is constrained by operator subjectivity. This study evaluates a Vision Transformer-based deep learning segmentation model for pancreatic tumors. Methods: A segmentation model using the USFM framework with a Vision Transformer backbone was trained and validated with 17,367 EUS images (from two public datasets) in 5-fold cross-validation. The model was tested on an independent dataset of 350 EUS images from another public dataset, manually segmented by radiologists. Preprocessing included grayscale conversion, cropping, and resizing to 512x512 pixels. Metrics included Dice similarity coefficient (DSC), intersection over union (IoU), sensitivity, specificity, and accuracy. Results: In 5-fold cross-validation, the model achieved a mean DSC of 0.651 +/- 0.738, IoU of 0.579 +/- 0.658, sensitivity of 69.8\%, specificity of 98.8\%, and accuracy of 97.5\%. For the external validation set, the model achieved a DSC of 0.657 (95\% CI: 0.634-0.769), IoU of 0.614 (95\% CI: 0.590-0.689), sensitivity of 71.8\%, and specificity of 97.7\%. Results were consistent, but 9.7\% of cases exhibited erroneous multiple predictions. Conclusions: The Vision Transformer-based model demonstrated strong performance for pancreatic tumor segmentation in EUS images. However, dataset heterogeneity and limited external validation highlight the need for further refinement, standardization, and prospective studies.}
}@misc{lopes2026alemdesempenhoumestudo,
      title={Al\'em do Desempenho: Um Estudo da Confiabilidade de Detectores de Deepfakes}, 
      author={Lucas Lopes and Rayson Laroca and André Grégio},
      year={2026},
      eprint={2601.08674},
      archivePrefix={arXiv},
      primaryClass={cs.CV},
      doi={https://doi.org/10.5753/sbseg.2025.11431},
      url={https://arxiv.org/abs/2601.08674},
      abstract = {Deepfakes are synthetic media generated by artificial intelligence, with positive applications in education and creativity, but also serious negative impacts such as fraud, misinformation, and privacy violations. Although detection techniques have advanced, comprehensive evaluation methods that go beyond classification performance remain lacking. This paper proposes a reliability assessment framework based on four pillars: transferability, robustness, interpretability, and computational efficiency. An analysis of five state-of-the-art methods revealed significant progress as well as critical limitations.}
}@misc{ruizfas2026zonebasedtrainingapproachlastmile,
      title={A zone-based training approach for last-mile routing using Graph Neural Networks and Pointer Networks}, 
      author={Àngel Ruiz-Fas and Carlos Granell and José Francisco Ramos and Joaquín Huerta and Sergio Trilles},
      year={2026},
      eprint={2601.04705},
      archivePrefix={arXiv},
      primaryClass={cs.LG},
      url={https://arxiv.org/abs/2601.04705},
      abstract = {Rapid e-commerce growth has pushed last-mile delivery networks to their limits, where small routing gains translate into lower costs, faster service, and fewer emissions. Classical heuristics struggle to adapt when travel times are highly asymmetric (e.g., one-way streets, congestion). A deep learning-based approach to the last-mile routing problem is presented to generate geographical zones composed of stop sequences to minimize last-mile delivery times.   The presented approach is an encoder-decoder architecture. Each route is represented as a complete directed graph whose nodes are stops and whose edge weights are asymmetric travel times. A Graph Neural Network encoder produces node embeddings that captures the spatial relationships between stops. A Pointer Network decoder then takes the embeddings and the route's start node to sequentially select the next stops, assigning a probability to each unvisited node as the next destination.   Cells of a Discrete Global Grid System which contain route stops in the training data are obtained and clustered to generate geographical zones of similar size in which the process of training and inference are divided. Subsequently, a different instance of the model is trained per zone only considering the stops of the training routes which are included in that zone.   This approach is evaluated using the Los Angeles routes from the 2021 Amazon Last Mile Routing Challenge. Results from general and zone-based training are compared, showing a reduction in the average predicted route length in the zone-based training compared to the general training. The performance improvement of the zone-based approach becomes more pronounced as the number of stops per route increases.}
}@misc{zhao2026attentionconsistencyregularizationinterpretable,
      title={Attention Consistency Regularization for Interpretable Early-Exit Neural Networks}, 
      author={Yanhua Zhao},
      year={2026},
      eprint={2601.08891},
      archivePrefix={arXiv},
      primaryClass={cs.LG},
      url={https://arxiv.org/abs/2601.08891},
      abstract = {Early-exit neural networks enable adaptive inference by allowing predictions at intermediate layers, reducing computational cost. However, early exits often lack interpretability and may focus on different features than deeper layers, limiting trust and explainability. This paper presents Explanation-Guided Training (EGT), a multi-objective framework that improves interpretability and consistency in early-exit networks through attention-based regularization. EGT introduces an attention consistency loss that aligns early-exit attention maps with the final exit. The framework jointly optimizes classification accuracy and attention consistency through a weighted combination of losses. Experiments on a real-world image classification dataset demonstrate that EGT achieves up to 98.97\% overall accuracy (matching baseline performance) with a 1.97x inference speedup through early exits, while improving attention consistency by up to 18.5\% compared to baseline models. The proposed method provides more interpretable and consistent explanations across all exit points, making early-exit networks more suitable for explainable AI applications in resource-constrained environments.}
}@misc{tlhomole2026operationalstreamflowforecastinglimpopo,
      title={Towards Operational Streamflow Forecasting in the Limpopo River Basin using Long Short-Term Memory Networks}, 
      author={James Tlhomole and Edoardo Borgomeo and Karthikeyan Matheswaran and Mariangel Garcia Andarcia},
      year={2026},
      eprint={2601.06941},
      archivePrefix={arXiv},
      primaryClass={cs.LG},
      url={https://arxiv.org/abs/2601.06941},
      abstract = {Robust hydrological simulation is key for sustainable development, water management strategies, and climate change adaptation. In recent years, deep learning methods have been demonstrated to outperform mechanistic models at the task of hydrological discharge simulation. Adoption of these methods has been catalysed by the proliferation of large sample hydrology datasets, consisting of the observed discharge and meteorological drivers, along with geological and topographical catchment descriptors. Deep learning methods infer rainfall-runoff characteristics that have been shown to generalise across catchments, benefitting from the data diversity in large datasets. Despite this, application to catchments in Africa has been limited. The lack of adoption of deep learning methodologies is primarily due to sparsity or lack of the spatiotemporal observational data required to enable downstream model training. We therefore investigate the application of deep learning models, including LSTMs, for hydrological discharge simulation in the transboundary Limpopo River basin, emphasising application to data scarce regions. We conduct a number of computational experiments primarily focused on assessing the impact of varying the LSTM model input data on performance. Results confirm that data constraints remain the largest obstacle to deep learning applications across African river basins. We further outline the impact of human influence on data-driven modelling which is a commonly overlooked aspect of data-driven large-sample hydrology approaches and investigate solutions for model adaptation under smaller datasets. Additionally, we include recommendations for future efforts towards seasonal hydrological discharge prediction and direct comparison or inclusion of SWAT model outputs, as well as architectural improvements.}
}@misc{pan2026disentanglinghistorypropagationdependencies,
      title={Disentangling History and Propagation Dependencies in Cross-Subject Knee Contact Stress Prediction Using a Shared MeshGraphNet Backbone}, 
      author={Zhengye Pan and Jianwei Zuo and Jiajia Luo},
      year={2026},
      eprint={2601.08318},
      archivePrefix={arXiv},
      primaryClass={q-bio.QM},
      url={https://arxiv.org/abs/2601.08318},
      abstract = {Background:Subject-specific finite element analysis accurately characterizes knee joint mechanics but is computationally expensive. Deep surrogate models provide a rapid alternative, yet their generalization across subjects under limited pose and load inputs remains unclear. It remains unclear whether the dominant source of prediction uncertainty arises from temporal history dependence or spatial propagation dependence. Methods:To disentangle these factors, we employed a shared MGN backbone with a fixed mesh topology. A dataset of running trials from nine subjects was constructed using an OpenSim-FEBio workflow. We developed four model variants to isolate specific dependencies: (1) a baseline MGN; (2) CT-MGN, incorporating a Control Transformer to encode short-horizon history; (3) MsgModMGN, applying state-conditioned modulation to message passing for adaptive propagation; (4) CT-MsgModMGN, combining both mechanisms. Models were evaluated using a rigorous grouped 3-fold cross-validation on unseen subjects.Results:The models incorporating history encoding significantly outperformed the baseline MGN and MsgModMGN in global accuracy and spatial consistency. Crucially, the CT module effectively mitigated the peak-shaving defect common in deep surrogates, significantly reducing peak stress prediction errors. In contrast, the spatial propagation modulation alone yielded no significant improvement over the baseline, and combining it with CT provided no additional benefit.Conclusion:Temporal history dependence, rather than spatial propagation modulation, is the primary driver of prediction uncertainty in cross-subject knee contact mechanics. Explicitly encoding short-horizon driver sequences enables the surrogate model to recover implicit phase information, thereby achieving superior fidelity in peak-stress capture and high-risk localization compared to purely state-based approaches.}
}
        --------------------
        